
GLUE was replaced by SuperGLUE within twelve months --- and no automated system detected the saturation that made this necessary. This work shows that such detection is feasible: by formalizing benchmark score-over-time trajectories as logistic growth processes and applying AIC model selection, decisive statistical justification for the logistic model is obtained ($\Delta AIC \approx -194$ to $-250$ versus linear, $-113$ to $-357$ versus power-law), and retrospective saturation detection is achieved within 3--5 months of community-recognized events. An empirically optimized criterion achieves 1-month accuracy.

The main contributions are: (1) to the authors' knowledge, the first automated benchmark saturation pipeline requiring no new test data collection; (2) formal information-theoretic validation of the logistic model for benchmark score dynamics; (3) a dual-criterion saturation detector with documented precision and baseline comparison; (4) discovery that logistic inflection points for GLUE and SuperGLUE predate their launch; and (5) preliminary boundary condition evidence extending the pipeline to $\geq$ 30 entries with bounded fitting.

One honest negative is reported: prospective forecasting from 6-month-early truncated data fails, scoping the method to retrospective analysis.

Future work should prioritize: HuggingFace API integration for live monitoring; prospective forecasting with shorter lookback windows; validation on real small benchmarks; and extension to vision and multilingual benchmarks.

If this pipeline had been running in 2019, it would have flagged GLUE's saturation within three months of the community-recognized event. The negative $t_0$ finding suggests the signal was present even before the leaderboard opened --- the benchmark was saturating before systematic measurement began. That observation, more than the detection accuracy itself, motivates continuing this line of work.

